\subsection{The exceptional orbital pair generates a six-dimensional unitary group}

\paragraph{D23.F25.}
Let $K_1,K_2,K_3$ be the literal Krammer matrices over
$\mathbb F_9=\mathbb F_3[i]$, $i^2=-1$, used in D23.P6--P9.
Set $X=K_1K_2$, $Y=K_3^2$, and $Z=(XY)^{-1}$.
There is a nondegenerate Hermitian form $H$ such that
\[
 \langle K_1,K_2,K_3\rangle=U(H),\qquad
 \langle X,Y\rangle=\det^{-1}(\{1,-1\})\subset U(H),
 \qquad \langle X,Y\rangle'=SU(H).
\tag{D23.F25.1}
\]
The respective group orders are
\[
 \begin{aligned}
 |U(H)|&=182699779456696320,\\
 |SU(H)|&=45674944864174080,\\
 |\langle X,Y\rangle|&=91349889728348160.
 \end{aligned}
\tag{D23.F25.2}
\]
Moreover $SU(H)$ is perfect, its centre is $\{I,-I\}$, and
$PSU(H)$ is simple.  The projective orders of $X,Y,Z$ are
respectively $36,6,3$, whereas their literal orders are $36,6,6$.
In particular the equal literal orders of $Y,Z$ do not give equal
conjugation-action orders on the derived group.

\begin{proof}
All matrix entries below use the encoding $a+3b\leftrightarrow a+bi$.
The exact invariant form is
\[
 H=\begin{pmatrix}
 0&1&1&6&6&0\\
 1&0&1&1&5&6\\
 1&1&0&0&1&1\\
 3&1&0&0&1&6\\
 3&8&1&1&0&1\\
 0&3&1&3&1&0
 \end{pmatrix}.
\tag{D23.F25.3}
\]
The certificate checks $H=\overline H^{\,T}$, $\operatorname{rank}H=6$,
and $\overline K_j^{\,T}HK_j=H$ for each $j$.  It also computes
the full invariant Hermitian-form space for $X,Y$, of dimension
one over $\mathbb F_3$.  The determinants are
\[
 \det K_j=i,\qquad \det X=\det Y=-1,\qquad\det Z=1.
\tag{D23.F25.4}
\]

There is a particularly small source of all the required root
elements.  The matrix $T=Y^2=K_3^4$ is
\[
 T=\begin{pmatrix}
 1&0&0&0&0&0\\0&1&0&0&0&0\\0&0&1&0&0&0\\
 0&0&0&1&0&0\\0&0&0&0&1&0\\0&2&3&2&3&1
 \end{pmatrix}.
\tag{D23.F25.5}
\]
Thus $\operatorname{rank}(T-I)=1$, $(T-I)^2=0$, and $T^3=I$.
Its image line is the sixth coordinate line.  An exact orbit
enumeration under $X,Y$ gives $22204$ distinct projective lines,
all isotropic for $H$.  This is the entire isotropic-line set,
not a sample.  To verify the count without a group-order table,
diagonalize a nondegenerate Hermitian form to the identity using
Gram--Schmidt and surjectivity of the norm $\mathbb F_9^\times
\to\mathbb F_3^\times$.  For a coordinate, the numbers of vectors
of norm $0,1,2$ are $1,4,4$.  If $a_n$ is the number of norm-zero
vectors in dimension $n$, then
\[
       a_1=1,\qquad a_n=4\cdot9^{n-1}-3a_{n-1}.
\]
This gives $a_6=177633$ and $(a_6-1)/8=22204$ isotropic lines.
The executable lists the complete orbit, with the first nonzero
coordinate normalized to one, and checks each norm directly.

We supply the generation argument needed to turn this orbit into
a group identification.  Write $h(v,w)=\overline v^{\,T}Hw$.
A rank-one unitary element $I+v\varphi$ with square-zero
$v\varphi$ fixes the hyperplane $\ker\varphi$ pointwise.
Unitarity implies that hyperplane is $v^\perp$; square-zero then
implies $h(v,v)=0$.  Consequently it has the form
\[
       T_{v,c}(w)=w+c\,h(v,w)v,\qquad \overline c=-c.
\tag{D23.F25.6}
\]
For one isotropic line over $\mathbb F_9$, there are exactly two
nonidentity such transformations, inverse to one another.
Conjugation carries their image line by the conjugating isometry.
Thus the orbit calculation proves that $\langle X,Y\rangle$
contains every unitary transvection of the six-dimensional space.

Here is an elementary proof that these transvections generate the
special unitary group in this situation.  On a hyperbolic plane
with form matrix $\left(\begin{smallmatrix}0&1\\1&0\end{smallmatrix}\right)$,
a determinant-one unitary matrix has diagonal entries in
$\mathbb F_3$ and off-diagonal entries in $i\mathbb F_3$.
Conjugating by $\operatorname{diag}(i,1)$ identifies this group
with $SL_2(\mathbb F_3)$.  Its upper and lower elementary matrices
are precisely the two transvection root groups, and elementary
row reduction generates $SL_2(\mathbb F_3)$ from them.
Every nondegenerate Hermitian plane over $\mathbb F_9$ is
hyperbolic: after diagonalization, norm surjectivity gives a
nonzero isotropic vector, and a hyperbolic partner is obtained
by subtracting half its norm times that vector.

The subgroup generated by all transvections is transitive on
vectors of norm one in every dimension $n\ge2$.  If two such
vectors span a nondegenerate plane, its special unitary group
carries one to the other: extend each to an orthonormal basis
of the plane and adjust the second vector by a norm-one scalar
to make determinant one.  This plane group is generated by
ambient transvections, as just proved.  Proportional vectors
also lie in such a plane.  In the remaining case write
$w=\alpha v+u$, where $h(v,v)=h(w,w)=1$, $u\ne0$ is isotropic
in $v^\perp$, and $N(\alpha)=1$.  Choose an isotropic partner
$u'$ in $v^\perp$ with $h(u,u')=1$.  Take $\gamma$ with
$N(\gamma)=2$, and choose $\beta$ so that
$\operatorname{Tr}(\overline\beta\gamma)=1$.
Then $z=\beta u+\gamma u'$ has norm one, is orthogonal to $v$,
and satisfies $N(h(w,z))=2$.  Both planes $\langle v,z\rangle$
and $\langle w,z\rangle$ are nondegenerate, giving the required
two-step move.  The trace choice exists because the nonzero
linear functional $\beta\mapsto\operatorname{Tr}
(\overline\beta\gamma)$ is onto $\mathbb F_3$.

Induction now proves generation of $SU_n$.  Given $g\in SU_n$,
first multiply by a product of transvections to fix a chosen
norm-one vector.  Its restriction to the orthogonal complement
belongs to $SU_{n-1}$ and is generated by its transvections by
induction.  The base $SU_1$ is trivial, and the plane case was
proved above.  This establishes the needed $SU_6$ generation
without invoking a classification of finite subgroups.

We have proved $SU(H)\subseteq\langle X,Y\rangle$.
Equation (D23.F25.4) proves the reverse containment in the
determinant preimage and supplies its nontrivial determinant
coset.  It also supplies all four determinant cosets in the
full braid image.  This proves the first two assertions of
(D23.F25.1).  Norm-one vector counts and stabilizers give
\[
 |U_n(3)|=3^{n(n-1)/2}\prod_{j=1}^{n}(3^j-(-1)^j),
 \qquad |SU_n(3)|=|U_n(3)|/4.
\]
For clarity, the stabilizer of a norm-one vector is $U_{n-1}$,
the group is transitive on those vectors by the argument above,
and their number is $3^{n-1}(3^n-(-1)^n)$, obtained from the
same norm recurrence.  This proves the order formula and
(D23.F25.2).

We also prove the structural facts used below.  In dimension
six, choose two orthogonal hyperbolic pairs, the first containing
the image vector of a given transvection.  An isometry which
scales that vector by $\alpha$ of norm two and its partner by
$\overline\alpha^{-1}$, and scales the second pair by
$\alpha^{-1},\overline\alpha$, has determinant one.
Its conjugation multiplies the transvection parameter by two.
Thus $[g,T]=T$ for the associated order-three transvection.
Every transvection is a commutator in $SU(H)$, so the generation
argument proves that $SU(H)$ is perfect.  Determinants now imply
that the derived groups of both groups in (D23.F25.1) are exactly
$SU(H)$.

To identify the centre and prove projective simplicity, use a
hyperbolic basis $e_1,e_2,e_3,f_1,f_2,f_3$.
An isometry fixing every isotropic line is diagonal on these
six lines.  The lines $e_i+e_j$, $f_i+f_j$ and $e_i+f_j$ for
$i\ne j$ then force all its diagonal entries to be equal.
It is therefore scalar.  A scalar in $SU(H)$ has fourth and
sixth power one, hence is $I$ or $-I$.  A central element
centralizes every transvection and therefore fixes each of
its image lines, proving that this is exactly the centre.

The action of $PSU(H)$ on isotropic lines is faithful and
primitive.  Here are explicit reasons for primitivity.  The
stabilizer of a line $v$ has exactly three orbits: $v$ itself,
the other lines orthogonal to $v$, and the lines not orthogonal
to $v$.  For nonorthogonal pairs, scale representatives to a
hyperbolic pair and complete to a hyperbolic basis.  For two
orthogonal independent lines, complete their span to two
hyperbolic pairs and then to a basis.  These constructions
give the asserted orbit transitivity; any determinant can be
corrected on the unused nondegenerate complement, without
moving the specified lines.  A block containing $v$ must be
a union of these stabilizer orbits.  The putative block of
all lines orthogonal to $v$ is excluded by moving $v$ to a
different orthogonal line $w$: the two perpendicular-line
sets intersect and differ.  The putative block consisting of
$v$ and its nonorthogonal lines is excluded by moving $v=e_1$
to $w=f_1$: these two sets intersect, but the isotropic line
$f_1+e_2$ belongs to the former and not the latter.
Thus only singleton and whole-set blocks remain.

For completeness, the elementary normal-subgroup argument is
as follows.  In a faithful primitive action, every nontrivial
normal subgroup is transitive, since its orbits are blocks.
Let $N$ be a nontrivial normal subgroup of $PSU(H)$ and let
$P$ be a point stabilizer.  Then $PSU(H)=NP$.  The order-three
transvection group at that point is abelian and normal in $P$;
its conjugates generate $PSU(H)$.  Modulo $N$, all these
conjugates have the same image, because every element is a
product from $N$ and $P$.  The quotient is consequently
abelian.  Perfectness forces it to be trivial, so $N=PSU(H)$.
This proves simplicity with no unproved finite-simple-group
identification.

Finally exact multiplication gives the local orders stated in
the assertion: no power $X^j$ with $0<j<36$ is scalar; no power
$Y^j$ with $0<j<6$ is scalar; and $Z^3=-I$ with $Z,Z^2$
nonscalar.  The centre of $U(H)$ is scalar by the same line
argument.  Moreover an element acting trivially by conjugation
on $PSU(H)$ acts trivially on $SU(H)$: its commutators would
otherwise define a homomorphism from the perfect group
$SU(H)$ to its centre.  It must therefore be scalar.
Thus these projective orders are also the conjugation-action
orders on both $SU(H)$ and $PSU(H)$.

The exact executable is
\path{agents/five_level_six/krammer_orbital_unitary.py}; its
JSON record is \path{krammer_orbital_unitary.json}, and its
complete line orbit is \path{krammer_orbital_isotropic_orbit.npz}
in the same directory.  The program checks only the stated
finite matrices and exhaustive line orbit.  Generation,
perfectness, simplicity and the subsequent Goursat deductions
are the mathematical proofs above, not assertions delegated
to a computer algebra group's name.
\end{proof}
